\documentclass[preview]{standalone}

\usepackage{amsmath}
\usepackage{amssymb}
\usepackage{stellar}
\usepackage{definitions}

\begin{document}

\id{prime-element-valuations}
\genpage

\section{Valuations associated with prime elements}

Let \(A\) be a \uniquefactorizationdomain, let
\[
    K=\quotientfield(A),
\]
and let \(p\in A\) be a \primeelement.

\begin{snippetdefinition}{prime-element-valuation-definition}{Valuation associated with a prime element}
    Let \(A\) be a \uniquefactorizationdomain, let
    \(K=\quotientfield(A)\), and let \(p\in A\) be a \primeelement.
    Every \(\alpha\in K^\times\) can be written as
    \[
        \alpha=u p_1^{n_1}\cdots p_r^{n_r},
    \]
    where \(u\in A^\times\), the \(p_i\) are pairwise nonassociate prime
    elements of \(A\), and \(n_i\in\integers\). The
    \emph{\(p\)-adic valuation} of \(\alpha\) is
    \[
        \nu_p(\alpha)
        =
        \begin{cases}
            n_i,&p\sim p_i,\\
            0,&p\text{ is not associate to any }p_i.
        \end{cases}
    \]
    Equivalently, if \(\alpha=a/b\) with
    \(a,b\in A\) and \(a,b\neq0_A\), then
    \[
        \nu_p(\alpha)=\nu_p(a)-\nu_p(b).
    \]
\end{snippetdefinition}

\begin{snippet}{prime-element-valuation-well-defined}
    Let \(A\) be a \uniquefactorizationdomain, let
    \(K=\quotientfield(A)\), and let \(p\in A\) be a \primeelement.
    For \(\alpha\in K^\times\), unique factorization shows that
    \(\nu_p(\alpha)\) does not depend on the chosen fraction
    \(\alpha=a/b\), nor on the representatives of the associate classes
    of the prime factors.
\end{snippet}

\begin{snippetexample}{rational-prime-valuations-example}{Valuations of a rational number}
    In \(A=\integers\), with \(K=\rationalnumbers\),
    \[
        \frac{27}{10}=2^{-1}3^3 5^{-1}.
    \]
    Therefore
    \[
        \nu_2\left(\frac{27}{10}\right)=-1,
        \qquad
        \nu_3\left(\frac{27}{10}\right)=3,
        \qquad
        \nu_5\left(\frac{27}{10}\right)=-1,
        \qquad
        \nu_7\left(\frac{27}{10}\right)=0.
    \]
\end{snippetexample}

\section{Valuation of a polynomial}

\begin{snippetdefinition}{prime-element-polynomial-valuation-definition}{Valuation of a polynomial}
    Let \(A\) be a \uniquefactorizationdomain, let
    \(K=\quotientfield(A)\), and let \(p\in A\) be a \primeelement. For a
    nonzero polynomial
    \[
        f=a_0+a_1x+\cdots+a_nx^n\in K[x],
    \]
    The \(p\)-adic valuation of \(f\) is
    \[
        \nu_p(f)=\min\{\nu_p(a_i)\suchthat a_i\neq0\}.
    \]
    By convention, \(\nu_p(0)=+\infty\).
\end{snippetdefinition}

\begin{snippetexample}{integer-polynomial-valuation-example}{Valuation of an integer polynomial}
    For
    \[
        f=2+6x+8x^2\in\integers[x],
    \]
    one has
    \[
        \nu_2(2)=1,
        \qquad
        \nu_2(6)=1,
        \qquad
        \nu_2(8)=3,
    \]
    and hence \(\nu_2(f)=1\).
\end{snippetexample}

\begin{snippetexample}{rational-polynomial-valuation-example}{Valuation of a rational polynomial}
    For
    \[
        f=5+\frac{1}{6}x+3x^2+\frac{7}{27}x^3
        \in\rationalnumbers[x],
    \]
    one has
    \[
        \nu_3(5)=0,
        \quad
        \nu_3\left(\frac{1}{6}\right)=-1,
        \quad
        \nu_3(3)=1,
        \quad
        \nu_3\left(\frac{7}{27}\right)=-3.
    \]
    Thus \(\nu_3(f)=-3\). Indeed,
    \[
        f
        =
        3^{-3}\left(135+\frac{9}{2}x+81x^2+7x^3\right),
    \]
    and the polynomial in parentheses has \(3\)-adic valuation \(0\).
\end{snippetexample}

\begin{snippetlemma}{polynomial-valuation-product-lemma}{Valuation of a product}
    Let \(A\) be a \uniquefactorizationdomain, let
    \(K=\quotientfield(A)\), and let \(p\in A\) be a
    \primeelement. For all nonzero \(f,g\in K[x]\),
    \[
        \nu_p(fg)=\nu_p(f)+\nu_p(g).
    \]
\end{snippetlemma}

\begin{snippetproof}{polynomial-valuation-product-lemma-proof}{polynomial-valuation-product-lemma}{Valuation of a product}
    First, for every \(\alpha\in K^\times\) and nonzero \(f\in K[x]\),
    \[
        \nu_p(\alpha f)=\nu_p(\alpha)+\nu_p(f),
    \]
    because multiplication by \(\alpha\) shifts the valuations of all
    coefficients by the same integer.

    We reduce to \(f,g\in A[x]\) with
    \(\nu_p(f)=\nu_p(g)=0\). Indeed, multiplying each polynomial by a
    suitable scalar in \(K^\times\) clears its coefficient
    denominators and removes its common \(p\)-power.

    Let
    \[
        \psi\colon A[x]\fromto(A/(p))[x]
    \]
    be coefficientwise reduction modulo \(p\). Since \(p\) is a
    \primeelement, \((p)\) is a prime \ideal, so \(A/(p)\) and hence
    \((A/(p))[x]\) are integral domains. The equalities
    \(\nu_p(f)=\nu_p(g)=0\) mean that \(\psi(f)\) and \(\psi(g)\) are
    nonzero. Therefore
    \[
        \psi(fg)=\psi(f)\psi(g)\neq0,
    \]
    which gives \(\nu_p(fg)=0\) in the reduced case.

    In general, choose \(\lambda,\mu\in K^\times\) such that
    \[
        \nu_p(\lambda)=-\nu_p(f),
        \qquad
        \nu_p(\mu)=-\nu_p(g),
    \]
    and \(\lambda f,\mu g\in A[x]\). Then
    \[
        \nu_p(\lambda f)=\nu_p(\mu g)=0.
    \]
    Applying the reduced case and the scalar identity gives
    \[
        0
        =
        \nu_p((\lambda f)(\mu g))
        =
        -\nu_p(f)-\nu_p(g)+\nu_p(fg),
    \]
    proving the claim.
\end{snippetproof}

\begin{snippetnote}{polynomial-reduction-kernel-valuation}{Kernel of coefficientwise reduction}
    Let \(A\) be a \uniquefactorizationdomain and let \(p\in A\) be a
    \primeelement. For coefficientwise reduction
    \[
        \psi\colon A[x]\fromto(A/(p))[x],
    \]
    one has
    \[
        \ringker\psi
        =
        \{h\in A[x]\suchthat p\text{ divides every coefficient of }h\}
        =
        \{h\in A[x]\suchthat \nu_p(h)\geq1\}.
    \]
\end{snippetnote}

\end{document}
